\documentclass[10pt,a4paper]{amsart}
\usepackage{amsmath,amssymb,amsthm,graphicx}
% Retained minimal preamble for the Stacks Project chapter more-algebra.tex (More on Algebra).
% The chapter file pinned-original-source.tex carries only \input{preamble} on its first line, so no
% preamble text lives inside the chapter itself.  The Stacks Project uses one single shared preamble for
% every chapter; the snapshot stages it identically as sources/stacks-chow/preamble.tex and
% sources/stacks-etale-cohomology/preamble.tex.  The block below is copied mechanically, line for line,
% from sources/stacks-chow/preamble.tex lines 14-18 (the Xy-pic setup) and 154-217 (the native theorem
% declarations together with the mathematical macros).  Nothing was added, removed or redefined; the
% wrapper class and the amsmath/amssymb/amsthm/graphicx base packages are supplied by the assembler.
% Chapter: more-algebra.tex -- More on Algebra.
% For commutative diagrams we use Xy-pic
\usepackage[all]{xy}

% We use 2cell for 2-commutative diagrams.
\xyoption{2cell}

% Theorem environments.
%
\theoremstyle{plain}
\newtheorem{theorem}[subsection]{Theorem}
\newtheorem{proposition}[subsection]{Proposition}
\newtheorem{lemma}[subsection]{Lemma}

\theoremstyle{definition}
\newtheorem{definition}[subsection]{Definition}
\newtheorem{example}[subsection]{Example}
\newtheorem{exercise}[subsection]{Exercise}
\newtheorem{situation}[subsection]{Situation}

\theoremstyle{remark}
\newtheorem{remark}[subsection]{Remark}
\newtheorem{remarks}[subsection]{Remarks}

\numberwithin{equation}{subsection}

% Macros
%
\def\lim{\mathop{\mathrm{lim}}\nolimits}
\def\colim{\mathop{\mathrm{colim}}\nolimits}
\def\Spec{\mathop{\mathrm{Spec}}}
\def\Hom{\mathop{\mathrm{Hom}}\nolimits}
\def\Ext{\mathop{\mathrm{Ext}}\nolimits}
\def\SheafHom{\mathop{\mathcal{H}\!\mathit{om}}\nolimits}
\def\SheafExt{\mathop{\mathcal{E}\!\mathit{xt}}\nolimits}
\def\Sch{\mathit{Sch}}
\def\Mor{\mathop{\mathrm{Mor}}\nolimits}
\def\Ob{\mathop{\mathrm{Ob}}\nolimits}
\def\Sh{\mathop{\mathit{Sh}}\nolimits}
\def\NL{\mathop{N\!L}\nolimits}
\def\CH{\mathop{\mathrm{CH}}\nolimits}
\def\proetale{{pro\text{-}\acute{e}tale}}
\def\etale{{\acute{e}tale}}
\def\QCoh{\mathit{QCoh}}
\def\Ker{\mathop{\mathrm{Ker}}}
\def\Im{\mathop{\mathrm{Im}}}
\def\Coker{\mathop{\mathrm{Coker}}}
\def\Coim{\mathop{\mathrm{Coim}}}

% Boxtimes
%
\DeclareMathSymbol{\boxtimes}{\mathbin}{AMSa}{"02}

%
% Macros for moduli stacks/spaces
%
\def\QCohstack{\mathcal{QC}\!\mathit{oh}}
\def\Cohstack{\mathcal{C}\!\mathit{oh}}
\def\Spacesstack{\mathcal{S}\!\mathit{paces}}
\def\Quotfunctor{\mathrm{Quot}}
\def\Hilbfunctor{\mathrm{Hilb}}
\def\Curvesstack{\mathcal{C}\!\mathit{urves}}
\def\Polarizedstack{\mathcal{P}\!\mathit{olarized}}
\def\Complexesstack{\mathcal{C}\!\mathit{omplexes}}
% \Pic is the operator that assigns to X its picard group, usage \Pic(X)
% \Picardstack_{X/B} denotes the Picard stack of X over B
% \Picardfunctor_{X/B} denotes the Picard functor of X over B
\def\Pic{\mathop{\mathrm{Pic}}\nolimits}
\def\Picardstack{\mathcal{P}\!\mathit{ic}}
\def\Picardfunctor{\mathrm{Pic}}
\def\Deformationcategory{\mathcal{D}\!\mathit{ef}}
\usepackage{cleveref}
\crefname{originallemma}{original lemma}{original lemmas}
\Crefname{originallemma}{Original lemma}{Original lemmas}
\makeatletter
\newlabel{more-algebra-lemma-open-mapping}{{0CQW}{source}{}{sourcepointer-1}{}}
\newlabel{more-algebra-lemma-open-mapping@cref}{{[originallemma][0][]0CQW}{[1][0][]source}{}{}{}}
\makeatother
\title{The open mapping lemma for complete linearly topologized groups}
\author{The Stacks Project Authors}
\date{stacks-project:more-algebra.tex:local-snapshot-20261012; GNU FDL 1.2 or later}
\begin{document}
\maketitle
\setcounter{subsection}{1}
\noindent\textit{Editorial note.} All delivered mathematical text is the author's own, copied from the pinned Stacks Project chapter file more-algebra.tex; no statement or proof step was rewritten, reordered or removed.\par\medskip
\noindent\textit{Editorial note.} The chapter file pinned-original-source.tex carries only \textbackslash{}input\{preamble\} on its first line, so no preamble text lives inside the chapter. The retained file sources/stacks-more-algebra/retained-preamble.tex is a mechanical, line-for-line copy of the project's single shared preamble as staged in the snapshot, namely sources/stacks-chow/preamble.tex lines 14-18 (Xy-pic setup) and 154-217 (theorem declarations and mathematical macros). Nothing was added, removed or redefined, no package was added, and the excerpt serves only as preamble, never as a span coordinate.\par\medskip
\noindent\textit{Editorial note.} Both counted spans start at the \textbackslash{}begin line of their environment, so the statement and the author's proof are delivered whole; surrounding prose of the section is deliberately outside the counted spans and is carried as standard prerequisites instead.\par\medskip
\noindent\textit{Editorial note.} The citation fields (source\_id, source\_version, authors, source\_url, license.name, source\_title) are identical to those of the cards already published from this chapter, so that the batch selector treats them as one source instead of deferring them as a same-source clash.\par\medskip
\noindent\textit{Allowed standard prerequisites (editorial).}\begin{itemize}
\item Linearly topologized abelian group with a countable fundamental system of neighbourhoods of zero, and completeness. the discussion of the same chapter, source lines 8996-9022; the neighbouring Baire category statement is counted by the sibling card S1-algebra-baire-category-complete-module at 9037-9075. Supplies the nested subgroups N\_n, the completeness used to sum the y\_i, and the openness of the closures of the images.
\item Nowhere dense subset and the failure of the second alternative. used at proof lines 9105-9106 to conclude that the closure of u(N\_n) is open. Is the negated hypothesis from which the whole construction starts.
\item Separatedness of the target and completeness of the source in the metric-like sense of Cauchy sequences in a filtered group. used at proof lines 9118-9119, where the infinite sum of the y\_i is formed and then identified with x\_1. Makes the element y exist and makes the equation x\_1 = u(y) valid, which is the crux of the proof.
\end{itemize}\medskip
\ifdefined\hypertarget\hypertarget{sourcepointer-1}{}\fi
\noindent\textit{Source reference.} Original lemma 0CQW: focused-seeds/AGP/sources/stacks-snapshot-20261011/tags.txt, source lines 15976--15976. The lemma's own canonical tag, quoted so that readers can resolve the label inside the delivered text; it points at the counted statement itself and carries no proof content.\par\medskip
\begin{lemma}[Open mapping lemma]
\label{lemma-open-mapping}
Let $u : N \to M$ be a continuous map of topologized
abelian groups. Assume that $M$ separated and that $N$ is complete, linearly
topologized, and has a countable fundamental system of neighbourhoods of $0$.
Then exactly one of the following holds
\begin{enumerate}
\item $u$ is open, or
\item for some open subgroup $N' \subset N$ the image
$u(N')$ is nowhere dense in $M$.
\end{enumerate}
\end{lemma}

\begin{proof}
Let $N_n$, $n \in \mathbf{N}$ be a fundamental system of neighbourhoods
of $0$. We may assume that $N_n$ is a subgroup and that
$N_{n + 1} \subset N_n$. If (2) does not hold, then
the closure $M_n$ of $u(N_n)$ is an open subgroup for $n = 1, 2, 3, \ldots$.
Since $u$ is continuous, we see that $M_n$, $n \in \mathbf{N}$
must be a fundamental system of open neighbourhoods of $0$ in $M$.
Also, since $M_n$ is the closure of $u(N_n)$ we see that
$$
u(N_n) + M_{n + 1} = M_n
$$
for all $n \geq 1$. Pick $x_1 \in M_1$. Then we can inductively choose
$y_i \in N_i$ and $x_{i + 1} \in M_{i + 1}$ such that
$$
u(y_i) + x_{i + 1} = x_i
$$
The element $y = y_1 + y_2 + y_3 + \ldots$ of $N$ exists because
$N$ is complete. Whereupon we see that $x_1 = u(y)$ because $M$ is separated.
Thus $M_1 = u(N_1)$. In exactly the same way the reader shows that
$M_i = u(N_i)$ for all $i \geq 2$ and we see that $u$ is open.
\end{proof}

\end{document}
